% Abhakseite „Das kann ich“ – Zone nur im Gesamt (\mitzone)
%% TODO Ich-kann-Sätze: Platzhalter wie in den Aufgabendateien
\begin{abhakseite}
\ifdefined\mitzone
\abhakgruppe{Kennst du schon}
\abhak{1}{Verschiebungen im ebenen Koordinatensystem ausführen und Punkte ablesen}
\abhak{2}{Satz des Pythagoras und Wurzelgleichungen (quadrieren, Lösungen sieben)}
\abhak{3}{Punkte im räumlichen Koordinatensystem lesen, Schrägbilder von Quadern und Prismen deuten}
\abhak{4}{Terme mit einem Buchstaben umformen und lineare Gleichungen lösen}
\abhak{5}{Mit Anteilen und Verhältnissen rechnen (ein Verhältnis mit einem Faktor ansetzen)}
\abhak{6}{Skalarprodukt in Koordinaten berechnen}
\abhak{7}{Satz des Pythagoras und Wurzelgleichungen (quadrieren, Lösungen sieben) – Fehler finden}
\abhak{8}{Satz des Pythagoras und Wurzelgleichungen (quadrieren, Lösungen sieben) – selbst rechnen}
\fi
\abhakgruppe{Einheit 1 · Vektorbegriff, Betrag und Kollinearität}
\abhak{9}{Faktor gesucht oder Richtung gesucht?}
\abhak{10}{Vektorbegriff und Betrag}
\abhak{11}{Vektorbegriff und Betrag – Fehler finden, Begründen, Darstellungswechsel, Anwendung}
\abhakgruppe{Einheit 2 · Vektorterme am Körper}
\abhak{12}{Vektorterme am Körper}
\abhak{13}{Vektorterme am Körper – Fehler finden, Begründen, Darstellungswechsel, Anwendung}
\abhakgruppe{Einheit 3 · Skalarprodukt im Sachzusammenhang}
\abhak{14}{Skalarprodukt im Sachzusammenhang}
\abhak{15}{Skalarprodukt im Sachzusammenhang – Fehler finden, Begründen, Anwendung}
\end{abhakseite}
